\chapter{Prototypische Umsetzung}


\begin{figure}[H]
\centering
\includegraphics[width=0.8\textwidth]{02_architektur.png}
\caption{Architektur der System-Pipeline}
\label{fig:architektur}
\end{figure}

\chapter{Verwendung BM25 und Exact-Match Signale}

Ein Bestandteil des vorgestellten Systems ist die Nutzung lexikalischer Retrievalsignale zur Beschreibung der Ähnlichkeit zwischen einer Anfrage und einem potentiellen EPLAN-Artikel \cite{EPLANArtikelverwaltung}.
Hier für wird unter anderem der BM25-Algorithmus verwendet, welcher ein etabliertes Verfahren im Bereich des Information Retrieval darstellt \cite{robertson2009prf}.
Anders als in klassischen Suchmaschinen, wird BM25 hier nicht als weitere Rankingstufe verwendet, sondern als eines von mehreren Signalen, die dem Learning-to-Rank-Modell als Eingabe dienen.
Die Aufgabe von diesem Algorithmus besteht deswegen darin, der Trainings-Kandidatenmenge ein erstes Relevanzsignal zu geben, um dadurch dem späteren Rankingmodell einen bewussten Bias in Richtung möglicher Query-Anfragen zu geben.


Neben BM25 wird auch ein Exact-Match-Algorithmus als Relevanzsignal verwendet. Dieser prüft die exakten Übereinstimmungen technischer Merkmale zwischen einer Datenblatt-Anfrage und den EPLAN-Artikeln \cite{EPLANArtikelverwaltung}.
Die Grundidee besteht darin, dass Artikel, die exakt die gleichen technischen Eigenschaften wie die Anfrage besitzen, eine höhere Relevanz aufweisen.

Das Vorliegende System erweitert dieses Konzept um eine gewichtete Bewertung technischer Eigenschaften.
Hierdurch wird eine differenzierte Relevanzbewertung ermöglicht, die die Anzahl und Bedeutung der übereinstimmenden Merkmale berücksichtigt.

Der Schwerpunkt liegt hierbei auf strukturierten technischen Merkmalen wie:

\begin{itemize}
  \item Herstellerkennungen
  \item Produktfamilien
  \item Produktcodes
  \item Makrostrukturen
  \item Kanalzahlen
  \item I/O-Schemata
  \item Positionsmerkmale
  \item Schutzarten
  \item Connector-Informationen
\end{itemize}



Nach der Berechnung von BM25 und Exact-Match liegen für jeden Trainingskandidaten (EPLAN-Artikel \cite{EPLANArtikelverwaltung}) zwei Relevanzsignale vor. 
BM25 bewertet die lexikalische Änlichkeit der Textuellen Repräsentationen. Exact-Match bewertet die spezifische Übereinstimmung technischer Merkmale.
Diese beiden Signale werden anschließend mittels Reciprocal Rank Fusion (RRF) zu einem gemeinsamen Retrievalsignal zusamengeführt \cite{cormack2009rrf}.




\subsection{Gewichtete RRF-Fusion im vorliegenden System} 

Im Gegensatz zur ursprünglichen Form verwendet die entwickelte Architektur eine gewichtete Variante von RRF. 
Der Grund dafür liegt in der unterschiedlichen Aussagekraft der beiden verwendeten Retrievalsignale. 
Während BM25 wertvolle Informationen über textuelle Ähnlichkeiten liefert, besitzen exakte technische Übereinstimmungen für die vorliegende Aufgabenstellung eine höhere fachliche Relevanz. 
Deswegen wird der Exact-Match-Zweig bewusst mit einer stärkeren Gewichtung als der BM25-Zweig versehen.

Die verwendete Fusionsfunktion lautet: 

\begin{align} 
    RRF(d) = 0.30 \cdot \frac{1} {30 + rank_{BM25}(d)} + 0.70 \cdot \frac{1} {30 + rank_{Exact}(d)} 
\end{align} 

mit 

\begin{description} 
    \item[$rank_{BM25}(d)$] Rangposition des Kandidaten in der BM25-Liste 
    \item[$rank_{Exact}(d)$] Rangposition des Kandidaten in der Exact-Match-Liste 
\end{description} 

Die Gewichtung von 0,70 für das Exact-Match-Signal und 0,30 für das BM25-Signal spiegelt die Priorisierung der technischen Übereinstimmungen wieder.
Dabei sollen Kandidaten mit identischen Produktcodes, Makrostrukturen oder technischen Eigenschaften stärker bevorzugt werden.
Kandidaten mit einer hohen textuellen Ähnlichkeit sollen entsprechend niedriger gewichtet werden.

Wichtig hierbei ist, dass RRF-Fusion nicht als eingenständige Rankingstufe verwendet wird, sondern lediglich als Aggregation der beiden Retrievalsignale BM25 und Exact-Match dient.
Dadurch entsteht ein weiteres Relevanzsignal, das dem Learning-to-Rank-Modell als Eingabe dient und die Stärken beider Verfahren kombiniert.